% GENERATED FILE. Edit canonical lessons, then run npm run generate:books.
% canonical-source-hash: fnv1a64:ffc06c4e46c7b8c8
% canonical-lessons: KA-C453-hadagu, KA-C453-navika, KA-C453-langaru, KA-C453-diksuci, KA-C453-sagara

\chapter{By Sea}
\label{ch:ka-by-sea}
\hlchaptermodality{453}

\begin{chapteropening}
\textbf{By the end of this chapter:} \emph{I can say a ship, a sailor, an anchor, a compass and an ocean.}

Complete the last lesson of chapter 453: I can say a ship, a sailor, an anchor, a compass and an ocean.
\end{chapteropening}

\section[\texorpdfstring{\kn{ಹಡಗು} (haḍagu)}{ಹಡಗು (haḍagu)}]{\kn{ಹಡಗು} (haḍagu) — a ship}

\label{lesson:KA-C453-hadagu}



\begin{quote}
Before the new one: say the Kannada for to transport, then the Kannada for a courier.
\end{quote}

\subsection*{You'll want to know: \kn{ಹಡಗು}}
\textbf{\kn{ಹಡಗು}} — \emph{haḍagu} — \textquotedblleft{}a ship\textquotedblright{}.

\begin{quote}
\kn{ಬಂದರಿನಲ್ಲಿ} \kn{ಒಂದು} \kn{ದೊಡ್ಡ} \kn{ಹಡಗು} \kn{ನಿಂತಿದೆ}.
\end{quote}

\begin{quote}
\emph{bandarinalli ondu doḍḍa haḍagu nintide.}
\end{quote}

\begin{quote}
A big ship is standing in the harbour.
\end{quote}

Plural: \kn{ಹಡಗುಗಳು}, \emph{haḍagugaḷu}.

A \kn{ದೋಣಿ} is a boat; a \kn{ಹಡಗು} crosses the sea.

\begin{cousinweb}
\textbf{\kn{ಹಡಗು}} (\emph{haḍagu}) is native Dravidian: DEDR 3838 lists Kannada \emph{paḍagu}, \emph{haḍagu}, the old \emph{p-} turned into \emph{h-}.
\end{cousinweb}

\begin{grammarlens}[title={What you've built}]
One more word for getting around.
\end{grammarlens}

\subsection*{Your turn}
\begin{itemize}
\raggedright
  \item \emph{Say it:} \emph{haḍagu}
  \item \emph{Say it:} \emph{haḍagu}, once more
  \item \emph{Say it:} \emph{bandarinalli ondu doḍḍa haḍagu nintide.} — A big ship is standing in the harbour.
  \item \emph{Recall:} say the Kannada for to transport, then the Kannada for a courier, then say \emph{haḍagu} again
\end{itemize}

\begin{tcolorbox}[breakable,colback=teal!4,colframe=teal!35!black,title={Before you move on}]
What does \kn{ಹಡಗು} mean? (\textquotedblleft{}A ship\textquotedblright{}.) Say it once more.
\end{tcolorbox}
\section[\texorpdfstring{\kn{ನಾವಿಕ} (nāvika)}{ನಾವಿಕ (nāvika)}]{\kn{ನಾವಿಕ} (nāvika) — a sailor}

\label{lesson:KA-C453-navika}



\begin{quote}
Before the new one: say the Kannada for a courier, then the Kannada for a ship.
\end{quote}

\subsection*{You'll want to know: \kn{ನಾವಿಕ}}
\textbf{\kn{ನಾವಿಕ}} — \emph{nāvika} — \textquotedblleft{}a sailor\textquotedblright{}.

\begin{quote}
\kn{ನನ್ನ} \kn{ಅಜ್ಜ} \kn{ಹಡಗಿನಲ್ಲಿ} \kn{ನಾವಿಕರಾಗಿದ್ದರು}.
\end{quote}

\begin{quote}
\emph{nanna ajja haḍaginalli nāvikarāgiddaru.}
\end{quote}

\begin{quote}
My grandfather was a sailor on a ship.
\end{quote}

Plural: \kn{ನಾವಿಕರು}, \emph{nāvikaru}.

Built on Sanskrit \emph{nau}, a boat, with -\kn{ಇಕ}, one who works with it: the one who works the boat.

\begin{cousinweb}
\textbf{\kn{ನಾವಿಕ}} (\emph{nāvika}) is borrowed from Sanskrit \emph{nāvika}, \textquotedblleft{}a sailor, a helmsman\textquotedblright{} (CDIAL 7082), from \emph{nau}, \textquotedblleft{}a boat\textquotedblright{}.
\end{cousinweb}

\begin{grammarlens}[title={What you've built}]
One more word for getting around.
\end{grammarlens}

\subsection*{Your turn}
\begin{itemize}
\raggedright
  \item \emph{Say it:} \emph{nāvika}
  \item \emph{Say it:} \emph{nāvika}, once more
  \item \emph{Say it:} \emph{nanna ajja haḍaginalli nāvikarāgiddaru.} — My grandfather was a sailor on a ship.
  \item \emph{Recall:} say the Kannada for a courier, then the Kannada for a ship, then say \emph{nāvika} again
\end{itemize}

\begin{tcolorbox}[breakable,colback=teal!4,colframe=teal!35!black,title={Before you move on}]
What does \kn{ನಾವಿಕ} mean? (\textquotedblleft{}A sailor\textquotedblright{}.) Say it once more.
\end{tcolorbox}
\section[\texorpdfstring{\kn{ಲಂಗರು} (laṅgaru)}{ಲಂಗರು (laṅgaru)}]{\kn{ಲಂಗರು} (laṅgaru) — an anchor}

\label{lesson:KA-C453-langaru}



\begin{quote}
Before the new one: say the Kannada for a ship, then the Kannada for a sailor.
\end{quote}

\subsection*{You'll want to know: \kn{ಲಂಗರು}}
\textbf{\kn{ಲಂಗರು}} — \emph{laṅgaru} — \textquotedblleft{}an anchor\textquotedblright{}.

\begin{quote}
\kn{ಹಡಗು} \kn{ಬಂದರಿನ} \kn{ಹೊರಗೆ} \kn{ಲಂಗರು} \kn{ಹಾಕಿತು}.
\end{quote}

\begin{quote}
\emph{haḍagu bandarina horage laṅgaru hākitu.}
\end{quote}

\begin{quote}
The ship dropped anchor outside the harbour.
\end{quote}

\kn{ಲಂಗರು} \kn{ಹಾಕು}, to drop anchor.

\begin{cousinweb}
\textbf{\kn{ಲಂಗರು}} (\emph{laṅgaru}) is usually traced to Persian \emph{langar}, \textquotedblleft{}an anchor\textquotedblright{}, which reached Kannada through Hindustani.
\end{cousinweb}

\begin{grammarlens}[title={What you've built}]
One more word for getting around.
\end{grammarlens}

\subsection*{Your turn}
\begin{itemize}
\raggedright
  \item \emph{Say it:} \emph{laṅgaru}
  \item \emph{Say it:} \emph{laṅgaru}, once more
  \item \emph{Say it:} \emph{haḍagu bandarina horage laṅgaru hākitu.} — The ship dropped anchor outside the harbour.
  \item \emph{Recall:} say the Kannada for a ship, then the Kannada for a sailor, then say \emph{laṅgaru} again
\end{itemize}

\begin{tcolorbox}[breakable,colback=teal!4,colframe=teal!35!black,title={Before you move on}]
What does \kn{ಲಂಗರು} mean? (\textquotedblleft{}An anchor\textquotedblright{}.) Say it once more.
\end{tcolorbox}
\section[\texorpdfstring{\kn{ದಿಕ್ಸೂಚಿ} (diksūci)}{ದಿಕ್ಸೂಚಿ (diksūci)}]{\kn{ದಿಕ್ಸೂಚಿ} (diksūci) — a compass}

\label{lesson:KA-C453-diksuci}



\begin{quote}
Before the new one: say the Kannada for a sailor, then the Kannada for an anchor.
\end{quote}

\subsection*{You'll want to know: \kn{ದಿಕ್ಸೂಚಿ}}
\textbf{\kn{ದಿಕ್ಸೂಚಿ}} — \emph{diksūci} — \textquotedblleft{}a compass\textquotedblright{}.

\begin{quote}
\kn{ದಿಕ್ಸೂಚಿಯ} \kn{ಮುಳ್ಳು} \kn{ಉತ್ತರಕ್ಕೆ} \kn{ತಿರುಗುತ್ತದೆ}.
\end{quote}

\begin{quote}
\emph{diksūciya muḷḷu uttarakke tiruguttade.}
\end{quote}

\begin{quote}
The needle of a compass turns to the north.
\end{quote}

\kn{ದಿಕ್ಕು}, a direction, plus \kn{ಸೂಚಿ}, a pointer.

\begin{cousinweb}
\textbf{\kn{ದಿಕ್ಸೂಚಿ}} (\emph{diksūci}) is built from two Sanskrit words: \emph{diś}, \textquotedblleft{}a direction\textquotedblright{} (CDIAL 6339), and \emph{sūci}, \textquotedblleft{}a needle\textquotedblright{} (CDIAL 13551).
\end{cousinweb}

\begin{grammarlens}[title={What you've built}]
One more word for getting around.
\end{grammarlens}

\subsection*{Your turn}
\begin{itemize}
\raggedright
  \item \emph{Say it:} \emph{diksūci}
  \item \emph{Say it:} \emph{diksūci}, once more
  \item \emph{Say it:} \emph{diksūciya muḷḷu uttarakke tiruguttade.} — The needle of a compass turns to the north.
  \item \emph{Recall:} say the Kannada for a sailor, then the Kannada for an anchor, then say \emph{diksūci} again
\end{itemize}

\begin{tcolorbox}[breakable,colback=teal!4,colframe=teal!35!black,title={Before you move on}]
What does \kn{ದಿಕ್ಸೂಚಿ} mean? (\textquotedblleft{}A compass\textquotedblright{}.) Say it once more.
\end{tcolorbox}
\section[\texorpdfstring{\kn{ಸಾಗರ} (sāgara)}{ಸಾಗರ (sāgara)}]{\kn{ಸಾಗರ} (sāgara) — an ocean}

\label{lesson:KA-C453-sagara}



\begin{quote}
Before the new one: say the Kannada for an anchor, then the Kannada for a compass.
\end{quote}

\subsection*{You'll want to know: \kn{ಸಾಗರ}}
\textbf{\kn{ಸಾಗರ}} — \emph{sāgara} — \textquotedblleft{}an ocean\textquotedblright{}.

\begin{quote}
\kn{ಹಡಗು} \kn{ಸಾಗರವನ್ನು} \kn{ಹತ್ತು} \kn{ದಿನದಲ್ಲಿ} \kn{ದಾಟಿತು}.
\end{quote}

\begin{quote}
\emph{haḍagu sāgaravannu hattu dinadalli dāṭitu.}
\end{quote}

\begin{quote}
The ship crossed the ocean in ten days.
\end{quote}

\kn{ಸಮುದ್ರ} is the sea; \kn{ಸಾಗರ} is the wider, grander word.

\begin{cousinweb}
\textbf{\kn{ಸಾಗರ}} (\emph{sāgara}) is borrowed from Sanskrit \emph{sāgara}, \textquotedblleft{}the ocean\textquotedblright{} (CDIAL 13325).
\end{cousinweb}

\begin{grammarlens}[title={What you've built}]
One more word for getting around.
\end{grammarlens}

\subsection*{Your turn}
\begin{itemize}
\raggedright
  \item \emph{Say it:} \emph{sāgara}
  \item \emph{Say it:} \emph{sāgara}, once more
  \item \emph{Say it:} \emph{haḍagu sāgaravannu hattu dinadalli dāṭitu.} — The ship crossed the ocean in ten days.
  \item \emph{Recall:} say the Kannada for an anchor, then the Kannada for a compass, then say \emph{sāgara} again
\end{itemize}

\begin{tcolorbox}[breakable,colback=teal!4,colframe=teal!35!black,title={Before you move on}]
What does \kn{ಸಾಗರ} mean? (\textquotedblleft{}An ocean\textquotedblright{}.) Say it once more.
\end{tcolorbox}
